\section{Extending an Entity You Do Not Own}
\label{sec:ch08-extending-a-foreign-entity}
\index{federation!extending a foreign entity|(}%

Mosaic still has a \code{Catalog} folder. It holds two files.

\begin{minted}{csharp}
[Key("id")]
public sealed class Product
{
    [ID]
    public required Guid Id { get; init; }

    [ReferenceResolver]
    public static Product? ResolveReference(string id, IResolverContext context)
        => ProductKey.TryDecode(id, context, out var productId)
            ? new Product { Id = productId }
            : null;
}
\end{minted}

That is the whole of Catalog inside Mosaic: a key, and how to build one of these
from a key. There is no lookup because there is nothing to look up. This service
does not know which products exist, cannot find out, and does not need to: it is
handed an identifier by something that does know, and its own four fields answer
against that identifier.

\index{federation!no @extends directive}%
Notice what is not here. There is no \code{@extends}, no \code{@external}, no
marker of any kind saying that another service owns this type. Under
Federation~v2 both subgraphs simply declare \code{Product} with the same key,
contribute whatever fields they have, and the composer works out that they are
the same type~\autocite{apollo2026entities}.
Chapter~\ref{ch:the-federation-model} said so in the abstract, and this is what
it looks like when it is your own code. The Mosaic side of \code{Product} is not
an extension of anything. It is a declaration that happens to be smaller.

\index{ID attribute@\code{[ID]} attribute!on a stub}%
The \code{[ID]} on the property rather than a resolver is deliberate and dull:
nothing in this service projects a product out of a database, so the reason
chapter~\ref{ch:data-without-the-n-plus-1} needed a resolver here does not
apply. What the attribute still has to do is encode, because what this field
answers is what a router will hand back to Catalog as a representation.

\subsection*{The published shape}

Here is what Mosaic now says about \code{Product}, from its own
\code{\_service}:

\begin{graphqlsdl}
type Product @key(fields: "id") {
  availableQuantity: Int!
  price: Money!
  reviews(first: Int, after: String, last: Int, before: String): ReviewConnection!
  averageRating: Float
  id: ID!
}
\end{graphqlsdl}

Four fields and a key, and \code{id} is last because it is the only one declared
on the class itself. The other four arrive from three
\code{[ObjectType<Product>]} partial classes in three different folders, in the
order the source generator found them.

\code{Product} does not implement \code{Node} here, and Catalog's does. That is
not an inconsistency to fix: a node resolver is a promise that an identifier is
enough to fetch the whole object, and Mosaic cannot keep that promise for a
product. It can answer four fields about one. Those are different claims and the
schema says so.

\subsection*{An order line hands out a key}

\index{federation!referencing a foreign entity}%
The other direction is where a service points at an entity it does not own.
\code{Order.lines} is Ordering's, and each line sold a product:

\begin{minted}{csharp}
public static Product GetProduct([Parent] OrderLine line)
    => new() { Id = line.ProductId };
\end{minted}

Until this chapter that method took a DataLoader, asked Catalog for the row, and
threw if it was missing. Now it wraps an identifier the line was already
storing. Ask Mosaic for an order, and this is all a line can say about what it
sold. Three things here are mine rather than the server's: the line breaks, the
removal of the customer's second order, and the removal of each line's
\code{unitPrice}, which is real and is not what this section is about.

\begin{minted}[fontsize=\footnotesize]{json}
{"data":{"ordersByCustomer":[
  {"id":"T3JkZXI6AAAA0AAAAECAAAAAAAAAAQ==",
   "total":{"amount":2812.00,"currency":"EUR"},
   "lines":[
     {"quantity":1,"product":{"id":"UHJvZHVjdDoAAACgAAAAQIAAAAAAAAAB"}},
     {"quantity":6,"product":{"id":"UHJvZHVjdDoAAACgAAAAQIAAAAAAAAAC"}},
     {"quantity":1,"product":{"id":"UHJvZHVjdDoAAACgAAAAQIAAAAAAAAAH"}}]}]}}
\end{minted}

The first line's \code{product.id} is character for character the string Catalog
answers for the same product. Two processes, no shared code, one identifier.
That is the entire contract, visible in one response, and it is the thing to
check first when a federated graph you inherited returns nulls where it should
return objects.

\index{correction!chapter 4's missing Include}%
That response could not have been produced at any earlier tag in this book, and
not for a federation reason. \code{OrderingService} never called
\code{.Include(o => o.Lines)}, and an order line is a related entity with a
shadow key rather than an owned type, so the navigation was never loaded:
\code{lines} answered with an empty array for every order and \code{total} threw
on the empty collection. Present since
chapter~\ref{ch:data-without-the-n-plus-1}, fixed here.

Chapter~\ref{ch:data-without-the-n-plus-1}'s prose is not wrong about any of
this; it says plainly that the lines are not owned, because an owned collection
is a navigation and \code{OrderLine} is a positional record. Only the service's
own comment claimed otherwise, and the missing call went with it. What is worth
carrying is why it lasted four chapters: no request in the verification
collection had ever asked an order for anything. This chapter found it by
accident, because \code{OrderLine.product} is the first place the book had a
reason to read an order line at all, and the gate now asks for one.

\medskip\noindent\textbf{What it costs.} The old version could tell you that a
line sold a product the catalog has never heard of, because it looked. This one
cannot, and neither can anything else in this service. A dangling identifier now
surfaces as a null product one hop away, in a process that will be asked about
it by a router rather than by the code that noticed.

\medskip\noindent\textbf{What it buys.} This resolver can no longer be slow. No
DataLoader, no batch, no statement. The order query got cheaper here and more
expensive somewhere else, and chapter~\ref{ch:enter-the-router} is where both
halves of that trade are visible at once.

\subsection*{The check that could not survive}

\index{federation!validation across a boundary}%
Chapter~\ref{ch:schema-design-that-survives-change} gave Mosaic a mutation with
four typed errors on its payload. \code{submitReview} now has three.

\code{ProductNotFoundError} is gone, because the check behind it is gone. That
resolver used to ask \code{CatalogService} whether the product existed before
writing the review, and the code comment beside it said, in
chapter~\ref{ch:hotchocolate-quickly}'s own words, that the check was on
borrowed time and that once Catalog was a separate service the honest answer
would be either a network call on the write path or an accepted risk.

This chapter takes the risk. An unknown product identifier is now stored without
complaint, and the union on the payload lost a member, which is a breaking
change for any client that was handling it. I want to be plain that I do not
love this. The alternative on the table was a synchronous call from Reviews to
Catalog on every write, which reintroduces exactly the coupling the extraction
was for and makes the write path fail when a service it does not depend on is
down. Chapter~\ref{ch:entity-resolution-done-right} has \code{@requires}, which
is as close as federation comes to giving the check back, and
chapter~\ref{ch:strangling-the-monolith} argues about whether it is worth
taking.

What I would not accept is leaving the error in the schema with nothing able to
raise it. An error type that can never occur is a lie told in a contract, and it
is worse than the missing check because a client will write a branch for it.

\index{federation!extending a foreign entity|)}%
